\section{The Limits of Human Intelligence}

My point is not to debate the merits of IQ tests, nor of what they measure. I recently discovered that my mother had saved all the school report cards, achievement test results, and an IQ test that I had taken at 10 years old. The IQ test is the motivation.

I can't find out much today about the \textbf{Kuhlmann-Anderson Test}, but it consisted of several verbal sections. Unlike other tests, there were no geometric pattern matching, and just one math section, so it was heavily weighted to verbal intelligence. There were two ways to get marked down:

\begin{enumerate}


\item Skip a problem due to time constraint

\item Provide the wrong answer
\end{enumerate}


Since the test was timed, there were usually a couple of uncompleted problems. So if the time was extended, perhaps as few as two or three minutes, I would have completed each section. Someone with a higher IQ than mine would have completed the test in the allotted time. That is why \textbf{Francis Galton} noticed a correlation between intelligence and reaction time.

So this leads to the first limit, at the individual level. In my case, with more time, I would have scored closer to those much more intelligent than I. However, I suspect that many others would not have performed much better because they still would have made too many mistakes. So their untimed score indicates the upper limit to their knowledge.

\paragraph{Academic Subjects}


As a student, I always preferred mathematics and hard sciences like physics and chemistry, to other subjects. That is because I never had to study them and had no problems with homework. Those subjects are considered the most difficult at universities. In my life I have tutored students in math, chemistry, and Latin, and had the misfortune to teach algebra to elementary school majors. That leads me to believe that there are hard limits to those students' abilities. Tutoring or teaching could only take them so far, and it is not simply a matter of lack of effort.

I was more successful teaching calculus; the students were self-selecting because they weren't forced to take the course. Even the poorer students could master basic calculus, even if they had to expend much more time and effort than the better students.

So those are the two examples:

\begin{itemize}


\item Students who are simply incapable of mastering a subject.

\item Students who can master the subject, given enough time and effort
\end{itemize}


\paragraph{CTMU}


\textbf{Christopher Langan}, who claims to be the most intelligent man

in the world, has developed a metaphysical system which he calls Cognitive Theoretic Model of the Universe (CTMU). His intelligence claim is based on having an alleged IQ of 195. Langan is prolific with many articles online and youtube videos.

From a high level, CTMU is interesting because the end result is a metaphysical system that looks a lot like Christianity without the dogmas, revelation, miracles, or exclusivity. There is God, the devil, an afterlife with heaven and hell, morality, and free will. There is also a political component which is basically extreme alt-right; he adds nothing new to political theory.

CTMU resolves all philosophical, religious, and metaphysical questions, if Langan can make a convincing case for it. So far, he has failed to do so. At a superficial level, some claims sound absurd, such as, ``God is the set of all sets'' and ``the universe is a set''. Neither makes sense, but he claims that, in the context of CTMU, they do make sense. He has promised to produce an axiomatic version of CTMU, but it has yet to appear; that could possibly silence his critics (which he accuses of incompetence).

These are the alternatives:

\begin{enumerate}


\item He can see CTMU faster than everyone else, who will eventually figure out CTMU.

\item He is so advanced that no other human can possibly understand it.
\end{enumerate}


For a summary, see: CTMU Wiki\footnote{\url{https://ctmucommunity.org/wiki/Main\_Page}}.

For a critique, see: \textit{Another crank comes to visit}\footnote{\url{http://www.goodmath.org/blog/2011/02/11/another-crank-comes-to-visit-the-cognitive-theoretic-model-of-the-universe/}}.

\paragraph{Hard Limits}


Are there hard limits to human understanding? If the world is intelligible, should not intelligent beings be able to understand it? Yet we see three issues:

\begin{enumerate}


\item Physics is at an impasse. Does it mean that fundamental scientific problems can never be solved?

\item Metaphysics: Philosophers have debated issues for millennia

\item Politics: No one can agree on the best regime
\end{enumerate}


Politics may well be the most difficult problem, or at least, the one that most urgently needs to be solved. What is the best way to organize a society? How to ensure that the most competent people are placed in positions of power?

The second point is the most baffling. In a democratic system, the masses are simply unable to recognize the most intelligent or most competent because they have no inner scale to make such a judgment. And worse, holding power requires a certain amount of ruthlessness or Machiavellianism, which may be antithetical to the social polity as a whole.

\paragraph{Chinese Civil Service Exam}


For over a thousand years, China relied on a civil service exam\footnote{\url{https://en.wikipedia.org/wiki/Imperial\_examination}} to select candidates to work in the state bureaucracy. It was open to everyone. The test included knowledge of Chinese culture and philosophy. This ensured social continuity and restricted civil service jobs to the most competent and intelligent.

That exam has been replaced in the People's Republic of China, with a test that evaluates IQ and general knowledge. See for sample questions\footnote{\url{https://www.straitstimes.com/asia/east-asia/chinas-civil-service-exam-can-you-answer-these-questions}}.

It remains to be seen if such a system will prevail.

\url{https://www.youtube.com/watch?v=QA0gjyXG5O0}

\flrightit{Posted on 2022-08-07 by Cologero}

\begin{center}* * *\end{center}

\begin{footnotesize}
\begin{sffamily}
\texttt{Mark Siegmund on 2022-08-08 at 10:29 said:}

thanks for an interesting piece... I did not know about Galton and intelligence/speed
\url{https://www.psychologicalscience.org/news/does-thinking-fast-mean-youre-thinking-smarter.html}

additionally, there are at least 3 EPGSIG members who have scored 200+, one of whom, is active on our Thousander FB site

\hfill

\texttt{Paulo Adolpho on 2022-08-11 at 21:23 said:}

Cologero, there are different types of intelligence, see psychologist Howard Gardner's theory of multiple intelligences. These IQ tests only analyze logical-mathematical intelligence. See the other forms of intelligence, such as musical, naturalistic, intrapersonal intelligence, etc.

\hfill

\texttt{m on 2022-08-15 at 11:12 said:}

``This ensured social continuity and restricted civil service jobs to the most competent and intelligent.''

Like any entrance test that is scored based upon `subjective' criteria or interpretation, the Chinese exam system was often abused, and open to bribery, political intrigue, and nepotism. Social continuity can be argued, however the most competent and intelligent? That was always an open question.

Compare Feng Menglong's vernacular works, Stories to Instruct the World, Stories to Caution the World, and Stories to Awaken the World... many examples of hapless scholars misused by the exam system, but in the end (usually) triumphed in spite of it. Alternately, stories of once worthy scholars who, after passing exams, became corrupted.

As an aside, I'd advise anyone unfamiliar with the author to approach his Wikipedia entry with a grain of salt. WikiP portrays the author as a kind of nascent `modern day' liberal, which he certainly was not. Nor were his stories based upon any kind of feminist orientation; especially the White Snake tale, that is now, in China (movies, film and more traditional opera), totally feminist.

\hfill

\texttt{Balrog Booger on 2022-08-16 at 13:00 said:}

The CTMU is very interesting, but I'm not competent to prove or disprove the mathematical soundness of it. He's addressed the notion of sets before. Regarding that, I'll include two paragraphs from his introduction to the system at the end of this comment. Assess his reasoning for yourselves.

I do think that Langan has become more wise since he gave the interview you linked. I don't know if he would disavow anything in particular that he said here, but his bearing has changed subtly but substantially these past few years.

His notion of a ``Human Singularity'' (as opposed to the technological singularity) is apt. It's a different symbolic expression of a worldwide return to Traditional society. He worked in earnest on the CTMU to try and assist in accomplishing that as much as he could, and if the theory is sound, he may have been a critical historical personage. Of course, that assumes that people would heed it if it was sound. We may be already be too corrupt to avert the Great Disaster either way.

I respect the man and appreciate his efforts in building this theory (as I esteem and thank Cologero for all the resources on this site). If there is some unsound part of the CTMU, the proper response is not to throw it out but to repair it. I think assisting in its construction, however we can with our individual abilities, is the solemn duty of our species. Langan is grappling with the Supreme object of knowledge. If you believe his effort is futile, so be it. But if you allow for any chance that such a theory could be built, you are morally bound to help that building. That is my opinion.

``Now for a brief word on sets. Mathematicians view set theory as fundamental. Anything can be considered an object, even a space or a process, and wherever there are objects, there is a set to contain them. This ``something'' may be a relation, a space or an algebraic system, but it is also a set; its relational, spatial or algebraic structure simply makes it a structured set. So mathematicians view sets, broadly including null, singleton, finite and infinite sets, as fundamental objects basic to meaningful descriptions of reality. It follows that reality itself should be a set... in fact, the largest set of all. But every set, even the largest one, has a powerset which contains it, and that which contains it must be larger (a contradiction). The obvious solution: define an extension of set theory incorporating two senses of ``containment'' which work together in such a way that the largest set can be defined as ``containing'' its powerset in one sense while being contained by its powerset in the other. Thus, it topologically includes itself in the act of descriptively including itself in the act of topologically including itself... , and so on, in the course of which it obviously becomes more than just a set.

``In the Cognitive-Theoretic Model of the Universe or CTMU, the set of all sets, and the real universe to which it corresponds, take the name (SCSPL) of the required extension of set theory. SCSPL, which stands for Self-Configuring Self-Processing Language, is just a totally intrinsic, i.e. completely self-contained, language that is comprehensively and coherently (self-distributively) self-descriptive, and can thus be model-theoretically identified as its own universe or referent domain. Theory and object go by the same name because unlike conventional ZF or NBG set theory, SCSPL hologically infuses sets and their elements with the distributed (syntactic, metalogical) component of the theoretical framework containing and governing them, namely SCSPL syntax itself, replacing ordinary set-theoretic objects with SCSPL syntactic operators. The CTMU is so-named because the SCSPL universe, like the set of all sets, distributively embodies the logical syntax of its own descriptive mathematical language. It is thus not only self-descriptive in nature; where logic denotes the rules of cognition (reasoning, inference), it is self-cognitive as well.''

\hfill

\texttt{Cologero on 2022-08-16 at 15:22 said:}

It is presumptive to assume that no one else in the history of the world has grappled with the Supreme object of knowledge. As I just pointed out, that Supreme object has revealed himself as well as the ``solemn duty'' of humanity. And you are arrogant to set my moral agenda since I have been writing about this Supreme object for 15 years.

Mathematicians view geometry as fundamental, but, alas, there are at least 3 geometries. Which one of them corresponds to reality? That is an empirical question, not a logical or tautological issue.

So if set theory is fundamental, then which one corresponds to reality: ZF, NBG, or CTMU? Until there is an axiomatic version of CTMU set theory, as was promised several years ago, we can't tell if it is consistent or not. As defined, ``containment'' is ambiguous as it is understood in two different senses.

What is an object? A set is a mathematical construct, but is it also an object? Apparently. Then so is every circle. Do you know how large the set of all circles is? Are physical things also object? But they are not consistent from moment to moment. How can there be a set whose members cannot be defined? These are just starting points.

There will be more tomorrow, based on a recent online ``debate''.

\hfill

\texttt{Balrog Booger on 2022-08-16 at 22:12 said:}

I said no such thing dismissing that no one else had grappled with the Supreme object of knowledge. I meant only to say that he is performing a work more weighty than the mundane physicist. I explicitly thanked you in the post for doing the same thing in your own way. I did not mean to accuse you of contradicting the task of awakening people to God. In this matter you have done much for me over the years. The fact that you criticize it from the perspective you do rules out true enmity toward it as a true atheist or a demon might have.

As I said in my previous comment, I'm not competent to address the soundness of the CTMU as defined by Langan, nor can I specify which within the three geometries you mention are fundamental. When I say that I think it is the solemn moral duty of people to assist with the CTMU, I shouldn't have written that comment so carelessly in this way. I don't see this matter as being something confined to the specific delineated theory of the CTMU, with all its baggage, but rather the matter in question is attaining undisputable specificity of metaphysical truth. The extent to which we can assist in the creation of the True, provable ``Theory of Everything,'' I believe we are morally bound to do so.

Langan, whatever his flaws, is trying his best to contribute to the true bettering of this world. I believe you are as well, honorable Cologero. Please forgive me if you feel I have been disrespectful or annoying.


\hfill

\end{sffamily}
\end{footnotesize}

